\section{Problem 50: derivatives of the totient distribution function (Round 7)}

\subsection{1) ROUND-7 OBJECTIVE}

\textbf{Path (C) (obstruction / correction).}
Round~6 proved a strong \emph{symmetric} residual oscillation theorem: on a dense $G_\delta$ set,
\[
\liminf_{h\to0^+}\frac{G(x+h)-G(x-h)}{2h}=0,
\qquad
\limsup_{h\to0^+}\frac{G(x+h)-G(x-h)}{2h}=+\infty.
\]
However, this still leaves open a genuine structural gap: the symmetric quotient does not distinguish the
left and right sides. In principle, a point could be flat on one side and wild on the other.

The new objective is therefore to strengthen Round~6 from \emph{symmetric} to \emph{one-sided} Dini derivatives.
We prove that there is a dense $G_\delta$ set of points (indeed a dense $G_\delta$ set disjoint from the special value set
$\mathcal E:=\{\varphi(n)/n:n\ge1\}$) at which \emph{both} one-sided lower Dini derivatives are $0$ and \emph{both}
one-sided upper Dini derivatives are $+\infty$.

Thus any hypothetical point where $G'(x)$ exists and is a finite positive number must avoid not only the dense countable set
$\mathcal E$ and the residual symmetric oscillation set from Round~6, but also a \emph{residual set of maximal one-sided oscillation}.

\subsection{2) Round-6 FOUNDATION USED}

We use the following previous-round results as black boxes.

\begin{itemize}
\item \textbf{Round~1.} The limiting distribution function
\[
G(u):=\lim_{X\to\infty}\frac1X\#\Bigl\{n\le X:\frac{\varphi(n)}n\le u\Bigr\}
\qquad(0\le u\le1)
\]
exists, is continuous and strictly increasing on $[0,1]$, and the associated measure $\mu=dG$ is purely singular with respect to Lebesgue measure.
In particular,
\[
G'(x)=0\qquad\text{for Lebesgue-a.e. }x\in(0,1).
\]

\item \textbf{Round~3.} The value set
\[
\mathcal E:=\Bigl\{\frac{\varphi(n)}{n}:n\ge1\Bigr\}
\]
is countable and dense in $[0,1]$, and for every $t\in\mathcal E$ one has
\[
\lim_{h\to0^+}\frac{G(t)-G(t-h)}{h}=+\infty.
\]
Equivalently, every $t\in\mathcal E$ has infinite \emph{left} upper Dini derivative.

\item \textbf{Round~6.} The symmetric flatness set
\[
A_{\rm sym}:=\Bigl\{x\in(0,1):\liminf_{h\to0^+}\frac{G(x+h)-G(x-h)}{2h}=0\Bigr\}
\]
is a dense $G_\delta$ subset of $(0,1)$ and has full Lebesgue measure.

\item \textbf{Round~6.} The symmetric blow-up set
\[
B_{\rm sym}:=\Bigl\{x\in(0,1):\limsup_{h\to0^+}\frac{G(x+h)-G(x-h)}{2h}=+\infty\Bigr\}
\]
is a dense $G_\delta$ subset of $(0,1)$.
\end{itemize}

\subsection{3) NEW INSIGHT / TOOL (ROUND-7)}

The new tool is a systematic passage from the symmetric quotient to the four one-sided Dini derivatives.
For $x\in(0,1)$ define
\[
\underline D_-G(x):=\liminf_{h\to0^+}\frac{G(x)-G(x-h)}{h},
\qquad
\overline D_-G(x):=\limsup_{h\to0^+}\frac{G(x)-G(x-h)}{h},
\]
\[
\underline D_+G(x):=\liminf_{h\to0^+}\frac{G(x+h)-G(x)}{h},
\qquad
\overline D_+G(x):=\limsup_{h\to0^+}\frac{G(x+h)-G(x)}{h}.
\]

The genuinely new theorem of this round is:
\begin{quote}
\emph{There exists a dense $G_\delta$ set $\Omega\subset(0,1)\setminus\mathcal E$ such that for every $x\in\Omega$,}
\[
\underline D_-G(x)=\underline D_+G(x)=0,
\qquad
\overline D_-G(x)=\overline D_+G(x)=+\infty.
\]
\end{quote}

This is strictly stronger than Round~6, which only gave the corresponding statement for the \emph{symmetric} quotient.

\subsection{4) ATTACK PLAN (ROUND-7)}

After Round~6, the exact gap was: generic symmetric oscillation does not force one-sided oscillation on \emph{each} side separately.
To close that gap we prove four statements.

\begin{enumerate}
\item The sets
\[
A_-:=\{x:\underline D_-G(x)=0\},
\qquad
A_+:=\{x:\underline D_+G(x)=0\}
\]
are both dense $G_\delta$ subsets of $(0,1)$, and in fact have full Lebesgue measure.

\item The sets
\[
B_-:=\{x:\overline D_-G(x)=+\infty\},
\qquad
B_+:=\{x:\overline D_+G(x)=+\infty\}
\]
are both dense $G_\delta$ subsets of $(0,1)$.

\item Therefore
\[
\Omega_0:=A_-\cap A_+\cap B_-\cap B_+
\]
is a dense $G_\delta$ subset of $(0,1)$.

\item Since $\mathcal E$ is countable, the refined set
\[
\Omega:=\Omega_0\setminus\mathcal E
\]
is still dense $G_\delta$.
\end{enumerate}

This overcomes the main Round~6 limitation: it shows that maximal oscillation occurs \emph{on both sides separately}, not merely after symmetrization.

\subsection{5) WORK (ROUND-7)}

\subsubsection*{5.1. One-sided quotients and a $G_\delta$ description}

For $0<h<1$ define the one-sided secant quotients
\[
Q_h^-(x):=\frac{G(x)-G(x-h)}{h}\qquad(h<x<1),
\]
\[
Q_h^+(x):=\frac{G(x+h)-G(x)}{h}\qquad(0<x<1-h).
\]
For fixed $h$, the functions $x\mapsto Q_h^-(x)$ and $x\mapsto Q_h^+(x)$ are continuous on their natural domains, since $G$ is continuous.
For fixed $x$, the maps $h\mapsto Q_h^-(x)$ and $h\mapsto Q_h^+(x)$ are also continuous for $0<h<\min\{x,1-x\}$.

For integers $N,m\ge1$ define
\[
V^-_{N,m}:=\bigcup_{q\in\mathbb Q\cap(0,1/m)}\Bigl\{x\in(q,1):Q_q^-(x)<1/N\Bigr\},
\]
\[
V^+_{N,m}:=\bigcup_{q\in\mathbb Q\cap(0,1/m)}\Bigl\{x\in(0,1-q):Q_q^+(x)<1/N\Bigr\},
\]
\[
U^-_{N,m}:=\bigcup_{q\in\mathbb Q\cap(0,1/m)}\Bigl\{x\in(q,1):Q_q^-(x)>N\Bigr\},
\]
\[
U^+_{N,m}:=\bigcup_{q\in\mathbb Q\cap(0,1/m)}\Bigl\{x\in(0,1-q):Q_q^+(x)>N\Bigr\}.
\]
Each of these sets is open in $(0,1)$.

\begin{lemma}\label{lem:Gdelta-onesided}
With the notation above,
\[
A_-:=\{x\in(0,1):\underline D_-G(x)=0\}=\bigcap_{N=1}^\infty\bigcap_{m=1}^\infty V^-_{N,m},
\]
\[
A_+:=\{x\in(0,1):\underline D_+G(x)=0\}=\bigcap_{N=1}^\infty\bigcap_{m=1}^\infty V^+_{N,m},
\]
\[
B_-:=\{x\in(0,1):\overline D_-G(x)=+\infty\}=\bigcap_{N=1}^\infty\bigcap_{m=1}^\infty U^-_{N,m},
\]
\[
B_+:=\{x\in(0,1):\overline D_+G(x)=+\infty\}=\bigcap_{N=1}^\infty\bigcap_{m=1}^\infty U^+_{N,m}.
\]
In particular, all four sets $A_-,A_+,B_-,B_+$ are $G_\delta$ subsets of $(0,1)$.
\end{lemma}

\begin{proof}
We prove the identity for $A_-$; the other three are identical.

Suppose first that $x\in A_-$. Then for every $N,m$ there exists $h\in(0,1/m)$ such that $Q_h^-(x)<1/N$.
By continuity of $h\mapsto Q_h^-(x)$, this strict inequality persists for all $h'$ in a neighborhood of $h$.
Choose a rational $q$ in that neighborhood with $0<q<1/m$. Then $Q_q^-(x)<1/N$, so $x\in V^-_{N,m}$.
Since $N,m$ were arbitrary, $x\in\bigcap_{N,m}V^-_{N,m}$.

Conversely, if $x\in\bigcap_{N,m}V^-_{N,m}$, then for every $N,m$ there exists a rational $q\in(0,1/m)$ with
$Q_q^-(x)<1/N$. Hence for every $N$ and every neighborhood of $0$ there exists $h$ in that neighborhood such that
$Q_h^-(x)<1/N$. This is exactly the statement that $\liminf_{h\to0^+}Q_h^-(x)=0$, i.e. $x\in A_-$.
\end{proof}

\subsubsection*{5.2. Residual and full-measure flatness on both sides}

\begin{theorem}[One-sided lower Dini derivatives vanish generically]\label{thm:lowers-zero}
The sets $A_-$ and $A_+$ are dense $G_\delta$ subsets of $(0,1)$.
Moreover, both have full Lebesgue measure.
\end{theorem}

\begin{proof}
By Lemma~\ref{lem:Gdelta-onesided}, both sets are $G_\delta$.
We next prove density.

Let $x\in A_{\rm sym}$, where $A_{\rm sym}$ is the dense $G_\delta$ set from Round~6.
Then there exists a sequence $h_n\to0^+$ such that
\[
\frac{G(x+h_n)-G(x-h_n)}{2h_n}\to0.
\]
But
\[
\frac{G(x+h_n)-G(x-h_n)}{2h_n}=\frac{Q_{h_n}^-(x)+Q_{h_n}^+(x)}{2},
\]
and both $Q_{h_n}^-(x)$ and $Q_{h_n}^+(x)$ are nonnegative because $G$ is increasing.
Therefore
\[
Q_{h_n}^-(x)\to0,
\qquad
Q_{h_n}^+(x)\to0.
\]
So $x\in A_-\cap A_+$. Hence
\[
A_{\rm sym}\subseteq A_-\cap A_+.
\]
Since $A_{\rm sym}$ is dense, both $A_-$ and $A_+$ are dense.

For the full-measure statement, Round~1 gives $G'(x)=0$ for Lebesgue-a.e. $x\in(0,1)$.
At every such point, both one-sided difference quotients converge to the same derivative $0$.
Hence $\underline D_-G(x)=\underline D_+G(x)=0$ for Lebesgue-a.e. $x$, so both $A_-$ and $A_+$ have full Lebesgue measure.
\end{proof}

\subsubsection*{5.3. Residual blow-up on both sides}

\begin{theorem}[One-sided upper Dini derivatives blow up on a residual set]\label{thm:uppers-infty}
The sets $B_-$ and $B_+$ are dense $G_\delta$ subsets of $(0,1)$.
\end{theorem}

\begin{proof}
By Lemma~\ref{lem:Gdelta-onesided}, both sets are $G_\delta$.

We first prove that $B_-$ is dense.
By Round~3, every $t\in\mathcal E$ satisfies
\[
\lim_{h\to0^+}\frac{G(t)-G(t-h)}{h}=+\infty,
\]
so $\mathcal E\subseteq B_-$. Since $\mathcal E$ is dense in $(0,1)$, the set $B_-$ is dense.

We now prove that $B_+$ is dense.
Fix a nonempty open interval $I=(a,b)\subset(0,1)$ and integers $N,m\ge1$.
Set
\[
\delta:=\frac{b-a}{3}>0.
\]
Because $\mathcal E$ is dense, choose
\[
t\in\mathcal E\cap(a+\delta,b-\delta).
\]
Again using Round~3, the left derivative at $t$ is $+\infty$, so there exists
\[
0<h<\min\{1/m,\delta/2\}
\]
with
\[
Q_h^-(t)=\frac{G(t)-G(t-h)}{h}>N.
\]
By continuity of $q\mapsto Q_q^-(t)$, the inequality $Q_q^-(t)>N$ also holds for all $q$ in some neighborhood of $h$.
Choose a rational
\[
q\in\mathbb Q\cap(0,1/m)
\]
in that neighborhood; in particular $q<\delta$. Set
\[
x:=t-q.
\]
Then
\[
a<t-\delta<x<t<b,
\]
so $x\in I$, and
\[
Q_q^+(x)=\frac{G(x+q)-G(x)}{q}=\frac{G(t)-G(t-q)}{q}=Q_q^-(t)>N.
\]
Hence $x\in I\cap U^+_{N,m}$.
Since $I,N,m$ were arbitrary, every $U^+_{N,m}$ is dense; therefore
\[
B_+=\bigcap_{N,m}U^+_{N,m}
\]
is dense.
\end{proof}

\subsubsection*{5.4. Main Round-7 theorem: residual maximal one-sided oscillation}

\begin{theorem}[Residual one-sided oscillation outside $\mathcal E$]\label{thm:main-round7}
The set
\[
\Omega_0:=A_-\cap A_+\cap B_-\cap B_+
\]
is a dense $G_\delta$ subset of $(0,1)$.
Consequently, the set
\[
\Omega:=\Omega_0\setminus\mathcal E
\]
is also a dense $G_\delta$ subset of $(0,1)$.
For every $x\in\Omega$ one has
\[
\underline D_-G(x)=\underline D_+G(x)=0,
\qquad
\overline D_-G(x)=\overline D_+G(x)=+\infty.
\]
In particular, no one-sided derivative of $G$ exists finitely at any point of $\Omega$, and therefore $G$ is not differentiable at any point of $\Omega$.
\end{theorem}

\begin{proof}
By Theorem~\ref{thm:lowers-zero}, the sets $A_-$ and $A_+$ are dense $G_\delta$.
By Theorem~\ref{thm:uppers-infty}, the sets $B_-$ and $B_+$ are dense $G_\delta$.
Therefore their intersection $\Omega_0$ is dense $G_\delta$.

Since $\mathcal E$ is countable, its complement $(0,1)\setminus\mathcal E$ is also a dense $G_\delta$ subset of $(0,1)$.
Thus
\[
\Omega=\Omega_0\cap\bigl((0,1)\setminus\mathcal E\bigr)
\]
is dense $G_\delta$ as well.
The displayed identities for the four Dini derivatives are exactly the defining properties of $A_-,A_+,B_-,B_+$.
If $G'(x)$ existed as a finite real number at some $x\in\Omega$, then all four one-sided Dini derivatives would equal that finite value, contradiction.
\end{proof}

\begin{corollary}[Sharper obstruction to a finite positive derivative]\label{cor:sharper-obstruction}
Let
\[
D_+:=\{x\in(0,1): G'(x)\text{ exists and }0<G'(x)<\infty\}.
\]
Then
\[
D_+\subseteq (0,1)\setminus\Omega.
\]
In particular, $D_+$ is meager, disjoint from $\mathcal E$, contained in a Lebesgue-null set, and disjoint from a dense $G_\delta$ set on which all four one-sided Dini derivatives are maximally oscillatory.
\end{corollary}

\begin{proof}
The inclusion $D_+\subseteq(0,1)\setminus\Omega$ follows immediately from Theorem~\ref{thm:main-round7}.
The meagerness follows because $\Omega$ is dense $G_\delta$.
Disjointness from $\mathcal E$ was already proved in Round~3.
Finally, by Round~1 we have $G'(x)=0$ for Lebesgue-a.e. $x$, so $D_+$ is contained in a Lebesgue-null set.
\end{proof}

\subsubsection*{5.5. What proof strategies are ruled out?}

Round~6 ruled out interval-wise regularity arguments based on the symmetric quotient.
Theorem~\ref{thm:main-round7} is stronger: on a residual set outside $\mathcal E$, \emph{each side separately} exhibits arbitrarily small and arbitrarily large secant slopes.
Thus any approach that first tries to obtain meaningful one-sided slope control on a nontrivial topological set cannot succeed.
In particular, no strategy that seeks differentiability by proving either
\[
\frac{G(x)-G(x-h)}{h}\asymp 1
\qquad\text{or}\qquad
\frac{G(x+h)-G(x)}{h}\asymp 1
\]
on a dense or residual set can be viable.

\subsection{6) ADVERSARIAL VERIFICATION}

\begin{itemize}
\item \textbf{Could symmetric flatness fail to imply one-sided flatness?}
No. If
\[
\frac{Q_h^-(x)+Q_h^+(x)}{2}\to0
\]
along a sequence and both terms are nonnegative, then each term must itself tend to $0$ along that sequence.
This is exactly why $A_{\rm sym}\subseteq A_-\cap A_+$.

\item \textbf{Is the $G_\delta$ description correct?}
Yes. The only subtle point is replacing an arbitrary scale $h$ by a rational scale $q$.
For fixed $x$, the maps $h\mapsto Q_h^\pm(x)$ are continuous on their natural domains, so strict inequalities survive under small perturbations of $h$.

\item \textbf{Does the proof of density of $B_+$ really use only established results?}
Yes. The proof uses only Round~3's fact that every $t\in\mathcal E$ has infinite \emph{left} derivative, together with density of $\mathcal E$ and continuity in the step-size parameter.
No new external theorem is invoked.

\item \textbf{Could $\Omega$ accidentally intersect $\mathcal E$?}
By construction it does not: we explicitly remove $\mathcal E$ at the final step.
This is important, because points of $\mathcal E$ already have a different certified pathology from Round~3.

\item \textbf{Endpoints.}
All one-sided quotients are considered only for $x\in(0,1)$ and sufficiently small $h$, so there is no boundary ambiguity.
\end{itemize}

\subsection{7) FINAL (EXACTLY ONE)}

\textbf{UNRESOLVED (BUT STRICTLY ADVANCED).}

Round~7 gives a genuine strengthening of Round~6: the residual oscillation phenomenon is now shown to occur \emph{on each side separately}.
There exists a dense $G_\delta$ set $\Omega\subset(0,1)\setminus\mathcal E$ such that for every $x\in\Omega$,
\[
\underline D_-G(x)=\underline D_+G(x)=0,
\qquad
\overline D_-G(x)=\overline D_+G(x)=+\infty.
\]
So, outside the special value set $\mathcal E$, there is still a residual set of points with maximal one-sided irregularity.

This strictly narrows the landscape for any hypothetical point where $G'(x)$ exists and is a finite positive number:
such a point must lie outside $\mathcal E$, outside a dense $G_\delta$ set of one-sided oscillation, and inside a Lebesgue-null exceptional set.
The original Erd\H{o}s question remains open, but another entire class of one-sided regularity strategies has now been ruled out.

\subsection{8) COMPLETION ESTIMATE (MANDATORY)}

\textbf{COMPLETION: 89\%}

\subsection{9) REFERENCES}

\begin{enumerate}
\item I. J. Schoenberg, \emph{On asymptotic distributions of arithmetical functions}, Trans. Amer. Math. Soc. \textbf{39} (1936), 315--330.
\item P. Erd\H{o}s, \emph{On the distribution function of additive functions}, Ann. of Math. (2) \textbf{47} (1946), 1--20.
\item G. Tenenbaum and V. Toulmonde, \emph{Sur le comportement local de la r\'epartition de l'indicatrice d'Euler}, Funct. Approx. Comment. Math. \textbf{35} (2006), 321--338.
\item V. Toulmonde, \emph{Module de continuit\'e de la fonction de r\'epartition de $\varphi(n)/n$}, Acta Arith. \textbf{121} (2006), no.~4, 367--402.
\end{enumerate}
